% ---------------------------------------------------------
% TABELA EMENTA 20: Química — Ano 1
% ---------------------------------------------------------
\begin{xltabular}{\linewidth}{|X|p{2.5cm}|p{2.5cm}|}
\hline
\multirow{3}{=}{\textcolor{ifscgreen}{\textbf{Unidade Curricular:}}\\[3pt]\textbf{\large Química — Ano 1}} & \multicolumn{2}{p{\dimexpr5cm+2\tabcolsep+\arrayrulewidth\relax}|}{\textcolor{ifscgreen}{\textbf{Semestre:}} \textbf{2º}} \\ \cline{2-3}
 & \textcolor{ifscgreen}{\textbf{CH EaD*:}} & \textcolor{ifscgreen}{\textbf{CH Total*:}} \\ \cline{2-3}
 & \textbf{00 h} & \textbf{40 h} \\ \hline
\endhead
\multicolumn{3}{|p{\dimexpr\linewidth-2\tabcolsep-2\arrayrulewidth\relax}|}{\textcolor{ifscgreen}{\textbf{Objetivos:}}} \\ \hline
\multicolumn{3}{|p{\dimexpr\linewidth-2\tabcolsep-2\arrayrulewidth\relax}|}{
\begin{itemize}[leftmargin=*,noitemsep,topsep=2pt]
 \item Estabelecer relações entre os conhecimentos sobre a matéria, sua constituição e suas transformações com situações do contexto social, tecnológico, ambiental e profissional.
 \item Compreender a Química como ciência, reconhecendo sua importância para a compreensão dos fenômenos naturais, para o desenvolvimento científico e tecnológico e para os processos produtivos.
 \item Identificar as propriedades da matéria, distinguindo substâncias puras e misturas e reconhecendo os principais métodos de separação empregados no cotidiano, na indústria e na preservação ambiental.
 \item Compreender a evolução dos modelos atômicos, reconhecendo sua importância para a construção do conhecimento científico.
 \item Identificar a estrutura do átomo, relacionando partículas subatômicas, distribuição eletrônica e propriedades dos elementos químicos.
 \item Interpretar a organização da Tabela Periódica, correlacionando a posição dos elementos às suas propriedades periódicas e aplicações.
 \item Compreender a formação das ligações químicas, relacionando os diferentes tipos de ligação às propriedades das substâncias.
 \item Reconhecer a geometria molecular e sua influência nas propriedades das substâncias, relacionando a polaridade das moléculas ao seu comportamento físico e químico.
 \item Identificar as forças intermoleculares, compreendendo sua influência nas propriedades macroscópicas dos materiais.
\end{itemize}
} \\ \hline
\multicolumn{3}{|p{\dimexpr\linewidth-2\tabcolsep-2\arrayrulewidth\relax}|}{\textcolor{ifscgreen}{\textbf{Conteúdos:}}} \\ \hline
\multicolumn{3}{|p{\dimexpr\linewidth-2\tabcolsep-2\arrayrulewidth\relax}|}{
\begin{itemize}[leftmargin=*,noitemsep,topsep=2pt]
 \item Química: objeto de estudo, aplicações e importância para a sociedade, o meio ambiente e os processos produtivos.
 \item Matéria: propriedades, estados físicos, substâncias puras e misturas.
 \item Processos de separação de misturas.
 \item Modelos atômicos.
 \item Estrutura atômica.
 \item Tabela periódica.
 \item Ligações químicas.
 \item Geometria molecular.
 \item Forças intermoleculares.
\end{itemize}
} \\ \hline
\multicolumn{3}{|p{\dimexpr\linewidth-2\tabcolsep-2\arrayrulewidth\relax}|}{\textcolor{ifscgreen}{\textbf{Estratégias de Ensino e Aprendizagem:}}} \\ \hline
\multicolumn{3}{|p{\dimexpr\linewidth-2\tabcolsep-2\arrayrulewidth\relax}|}{- **Metodologia:** Metodologias diversificadas, priorizando a contextualização dos conceitos químicos em situações do cotidiano, dos processos produtivos e das questões ambientais. Aulas expositivas dialogadas com valorização dos conhecimentos prévios; resolução de exercícios; estudos dirigidos; modelos didáticos e representações da estrutura da matéria; jogos didáticos; vídeos; softwares de simulação e animação; atividades experimentais no Laboratório de Química (propriedades da matéria, separação de misturas, identificação de substâncias). SIGAA utilizado como Ambiente Virtual de Aprendizagem. Nas atividades práticas de laboratório, as turmas poderão ser divididas em subgrupos.
- **Avaliação:** Contínua e processual, com relatórios de práticas laboratoriais, listas de exercícios, estudos dirigidos, seminários, participação nas atividades propostas, produções individuais e coletivas, e avaliações individuais escritas.} \\ \hline
\multicolumn{3}{|p{\dimexpr\linewidth-2\tabcolsep-2\arrayrulewidth\relax}|}{\textcolor{ifscgreen}{\textbf{Bibliografia Básica:}}} \\ \hline
\multicolumn{3}{|p{\dimexpr\linewidth-2\tabcolsep-2\arrayrulewidth\relax}|}{1. CHANG, Raymond. \textbf{Química geral: conceitos essenciais}. 4. ed. Porto Alegre: AMGH, 2010.
2. FELTRE, Ricardo. \textbf{Química 1: química geral}. 7. ed. São Paulo: Moderna, 2008.
3. FRANCO, Dalton. \textbf{360º: química: cotidiano e transformações: volume único}. São Paulo: FTD, 2015.
4. PERUZZO, Francisco Miragaia; CANTO, Eduardo Leite do. \textbf{Química na abordagem do cotidiano: volume único}. 3. ed. São Paulo: Moderna, 2007.} \\ \hline
\multicolumn{3}{|p{\dimexpr\linewidth-2\tabcolsep-2\arrayrulewidth\relax}|}{\textcolor{ifscgreen}{\textbf{Bibliografia Complementar:}}} \\ \hline
\multicolumn{3}{|p{\dimexpr\linewidth-2\tabcolsep-2\arrayrulewidth\relax}|}{1. ATKINS, P. W.; JONES, Loretta. \textbf{Princípios de química: questionando a vida moderna e o meio ambiente}. 5. ed. Porto Alegre: Bookman, 2012.
2. BRADY, James E.; HUMISTON, Gerard E. \textbf{Química geral}. 2. ed. Rio de Janeiro: LTC, 2011. v. 1.
3. USBERCO, João; SALVADOR, Edgard. \textbf{Química: volume único}. 5. ed. reform. São Paulo: Saraiva, 2002.} \\ \hline
\end{xltabular}

\vspace{0.4cm}


% ---------------------------------------------------------
% TABELA EMENTA 21: Química — Ano 2
% ---------------------------------------------------------
\begin{xltabular}{\linewidth}{|X|p{2.5cm}|p{2.5cm}|}
\hline
\multirow{3}{=}{\textcolor{ifscgreen}{\textbf{Unidade Curricular:}}\\[3pt]\textbf{\large Química — Ano 2}} & \multicolumn{2}{p{\dimexpr5cm+2\tabcolsep+\arrayrulewidth\relax}|}{\textcolor{ifscgreen}{\textbf{Semestre:}} \textbf{2º e 3º}} \\ \cline{2-3}
 & \textcolor{ifscgreen}{\textbf{CH EaD*:}} & \textcolor{ifscgreen}{\textbf{CH Total*:}} \\ \cline{2-3}
 & \textbf{00 h} & \textbf{80 h} \\ \hline
\endhead
\multicolumn{3}{|p{\dimexpr\linewidth-2\tabcolsep-2\arrayrulewidth\relax}|}{\textcolor{ifscgreen}{\textbf{Objetivos:}}} \\ \hline
\multicolumn{3}{|p{\dimexpr\linewidth-2\tabcolsep-2\arrayrulewidth\relax}|}{
\begin{itemize}[leftmargin=*,noitemsep,topsep=2pt]
 \item Estabelecer relações entre os conhecimentos sobre as transformações químicas e situações do contexto social, tecnológico, ambiental e profissional.
 \item Identificar as principais funções inorgânicas, reconhecendo suas propriedades, nomenclatura, aplicações e ocorrência no cotidiano, nos processos produtivos e no meio ambiente.
 \item Compreender as reações químicas como processos de transformação da matéria, representando-as por meio de equações químicas corretamente balanceadas.
 \item Aplicar as relações estequiométricas na resolução de problemas envolvendo reagentes, produtos e rendimento das reações químicas.
 \item Compreender os conceitos relacionados às soluções, interpretando diferentes formas de expressar concentração.
 \item Analisar as transformações energéticas envolvidas nas reações químicas, compreendendo os princípios da termoquímica.
 \item Compreender os fatores que influenciam a velocidade das reações químicas e o estabelecimento do equilíbrio químico.
 \item Compreender os princípios da eletroquímica, relacionando as reações de oxirredução ao funcionamento de pilhas, baterias e processos eletrolíticos.
 \item Reconhecer os fundamentos da radioatividade, identificando suas aplicações, benefícios, riscos e implicações para a saúde, o meio ambiente e o desenvolvimento científico e tecnológico.
\end{itemize}
} \\ \hline
\multicolumn{3}{|p{\dimexpr\linewidth-2\tabcolsep-2\arrayrulewidth\relax}|}{\textcolor{ifscgreen}{\textbf{Conteúdos:}}} \\ \hline
\multicolumn{3}{|p{\dimexpr\linewidth-2\tabcolsep-2\arrayrulewidth\relax}|}{
\begin{itemize}[leftmargin=*,noitemsep,topsep=2pt]
 \item Funções Inorgânicas.
 \item Reações químicas e balanceamento de equações.
 \item Relações estequiométricas.
 \item Soluções.
 \item Termoquímica.
 \item Cinética química.
 \item Equilíbrio químico.
 \item Eletroquímica.
 \item Radioatividade.
\end{itemize}
} \\ \hline
\multicolumn{3}{|p{\dimexpr\linewidth-2\tabcolsep-2\arrayrulewidth\relax}|}{\textcolor{ifscgreen}{\textbf{Estratégias de Ensino e Aprendizagem:}}} \\ \hline
\multicolumn{3}{|p{\dimexpr\linewidth-2\tabcolsep-2\arrayrulewidth\relax}|}{- **Metodologia:** Metodologias diversificadas, compatíveis com os objetivos de aprendizagem, priorizando a contextualização dos conceitos químicos em situações do cotidiano, dos processos produtivos, da saúde e do meio ambiente. Aulas expositivas dialogadas, resolução de exercícios, estudos dirigidos, atividades investigativas, experimentação em laboratório, recursos audiovisuais e digitais, simulações computacionais, jogos didáticos, seminários e atividades individuais e em grupo. SIGAA utilizado como Ambiente Virtual de Aprendizagem. Nas atividades práticas de laboratório, as turmas poderão ser divididas em subgrupos.
- **Avaliação:** Contínua e processual, por meio de atividades escritas, relatórios, listas de exercícios, seminários, práticas laboratoriais, produções individuais e coletivas e avaliações individuais.} \\ \hline
\multicolumn{3}{|p{\dimexpr\linewidth-2\tabcolsep-2\arrayrulewidth\relax}|}{\textcolor{ifscgreen}{\textbf{Bibliografia Básica:}}} \\ \hline
\multicolumn{3}{|p{\dimexpr\linewidth-2\tabcolsep-2\arrayrulewidth\relax}|}{1. CHANG, Raymond. \textbf{Química geral: conceitos essenciais}. 4. ed. Porto Alegre: AMGH, 2010.
2. ATKINS, P. W.; JONES, Loretta. \textbf{Princípios de química: questionando a vida moderna e o meio ambiente}. 5. ed. Porto Alegre: Bookman, 2012.
3. PERUZZO, Francisco Miragaia; CANTO, Eduardo Leite do. \textbf{Química na abordagem do cotidiano: volume único}. 3. ed. São Paulo: Moderna, 2007.} \\ \hline
\multicolumn{3}{|p{\dimexpr\linewidth-2\tabcolsep-2\arrayrulewidth\relax}|}{\textcolor{ifscgreen}{\textbf{Bibliografia Complementar:}}} \\ \hline
\multicolumn{3}{|p{\dimexpr\linewidth-2\tabcolsep-2\arrayrulewidth\relax}|}{1. BRADY, James E.; HUMISTON, Gerard E. \textbf{Química geral}. 2. ed. Rio de Janeiro: LTC, 2014. v. 2.
2. FELTRE, Ricardo. \textbf{Química: físico-química, volume 2}. 7. ed. São Paulo: Moderna, 2008.
3. FRANCO, Dalton. \textbf{360º: química: cotidiano e transformações: volume único}. São Paulo: FTD, 2015.} \\ \hline
\end{xltabular}

\vspace{0.4cm}


% ---------------------------------------------------------
% TABELA EMENTA 22: Química — Ano 3
% ---------------------------------------------------------
\begin{xltabular}{\linewidth}{|X|p{2.5cm}|p{2.5cm}|}
\hline
\multirow{3}{=}{\textcolor{ifscgreen}{\textbf{Unidade Curricular:}}\\[3pt]\textbf{\large Química — Ano 3}} & \multicolumn{2}{p{\dimexpr5cm+2\tabcolsep+\arrayrulewidth\relax}|}{\textcolor{ifscgreen}{\textbf{Semestre:}} \textbf{1º (Química III)}} \\ \cline{2-3}
 & \textcolor{ifscgreen}{\textbf{CH EaD*:}} & \textcolor{ifscgreen}{\textbf{CH Total*:}} \\ \cline{2-3}
 & \textbf{00 h} & \textbf{40 h} \\ \hline
\endhead
\multicolumn{3}{|p{\dimexpr\linewidth-2\tabcolsep-2\arrayrulewidth\relax}|}{\textcolor{ifscgreen}{\textbf{Objetivos:}}} \\ \hline
\multicolumn{3}{|p{\dimexpr\linewidth-2\tabcolsep-2\arrayrulewidth\relax}|}{
\begin{itemize}[leftmargin=*,noitemsep,topsep=2pt]
 \item Estabelecer relações entre os conhecimentos de Química Orgânica e situações do contexto social, ambiental e profissional.
 \item Identificar as propriedades do átomo de carbono que justificam a diversidade dos compostos orgânicos e sua importância para os seres vivos, os materiais e os processos industriais.
 \item Classificar as cadeias carbônicas quanto à sua estrutura, correlacionando-as às propriedades dos compostos orgânicos.
 \item Reconhecer os principais grupos funcionais da Química Orgânica, relacionando suas propriedades físicas e químicas às suas aplicações no cotidiano.
 \item Compreender os conceitos de isomeria plana e espacial, identificando como diferenças estruturais podem influenciar as propriedades e aplicações dos compostos orgânicos.
 \item Conhecer as principais reações orgânicas, reconhecendo sua relevância em processos naturais, industriais e tecnológicos.
\end{itemize}
} \\ \hline
\multicolumn{3}{|p{\dimexpr\linewidth-2\tabcolsep-2\arrayrulewidth\relax}|}{\textcolor{ifscgreen}{\textbf{Conteúdos:}}} \\ \hline
\multicolumn{3}{|p{\dimexpr\linewidth-2\tabcolsep-2\arrayrulewidth\relax}|}{
\begin{itemize}[leftmargin=*,noitemsep,topsep=2pt]
 \item O átomo de carbono e suas propriedades.
 \item Cadeias carbônicas: definição e classificações.
 \item Principais grupos funcionais e suas propriedades (hidrocarbonetos, álcool, enol, fenol, éter, aldeído, cetona, ácido carboxílico, éster, amina e amida).
 \item Isomeria Plana e Espacial.
 \item Noções de Reações Orgânicas.
 \item Temáticas transversais: meio ambiente, saúde, educação alimentar e nutricional.
\end{itemize}
} \\ \hline
\multicolumn{3}{|p{\dimexpr\linewidth-2\tabcolsep-2\arrayrulewidth\relax}|}{\textcolor{ifscgreen}{\textbf{Estratégias de Ensino e Aprendizagem:}}} \\ \hline
\multicolumn{3}{|p{\dimexpr\linewidth-2\tabcolsep-2\arrayrulewidth\relax}|}{- **Metodologia:** Metodologias diversificadas, priorizando situações reais vinculadas ao cotidiano dos estudantes. Aulas expositivas dialogadas, com valorização dos conhecimentos prévios; resolução de exercícios; estudos dirigidos; jogos didáticos; vídeos e filmes; seminários; atividades experimentais no Laboratório de Química. Recursos: quadro, projetor multimídia, recursos audiovisuais, softwares de simulação e animação, materiais e equipamentos laboratoriais. SIGAA utilizado como Ambiente Virtual de Aprendizagem. Nas atividades práticas de laboratório, as turmas poderão ser divididas em subgrupos.
- **Avaliação:** Contínua e processual, com relatórios de práticas laboratoriais, listas de exercícios, estudos dirigidos, seminários, participação nas atividades propostas, produções individuais e coletivas, e avaliações individuais escritas.} \\ \hline
\multicolumn{3}{|p{\dimexpr\linewidth-2\tabcolsep-2\arrayrulewidth\relax}|}{\textcolor{ifscgreen}{\textbf{Bibliografia Básica:}}} \\ \hline
\multicolumn{3}{|p{\dimexpr\linewidth-2\tabcolsep-2\arrayrulewidth\relax}|}{1. ATKINS, P. W.; JONES, Loretta. \textbf{Princípios de química: questionando a vida moderna e o meio ambiente}. 5. ed. Porto Alegre: Bookman, 2012.
2. MCMURRY, John. \textbf{Química orgânica: combo}. 3. ed. São Paulo: Cengage Learning, 2016.
3. Livro didático fornecido pelo Fundo Nacional de Desenvolvimento da Educação (FNDE).} \\ \hline
\multicolumn{3}{|p{\dimexpr\linewidth-2\tabcolsep-2\arrayrulewidth\relax}|}{\textcolor{ifscgreen}{\textbf{Bibliografia Complementar:}}} \\ \hline
\multicolumn{3}{|p{\dimexpr\linewidth-2\tabcolsep-2\arrayrulewidth\relax}|}{1. CHANG, Raymond. \textbf{Química geral: conceitos essenciais}. 4. ed. Porto Alegre: AMGH, 2010.
2. DALMAZ, Carla; CALCAGNOTTO, Maria Elisa (revisão técnica). \textbf{Bioquímica ilustrada}. 7. ed. Porto Alegre: Artmed, 2019.
3. ZUBRICK, James W. \textbf{Manual de sobrevivência no laboratório de química orgânica: guia de técnicas para o aluno}. Rio de Janeiro: LTC, 2005.} \\ \hline
\end{xltabular}

\vspace{0.4cm}


% ---------------------------------------------------------
% TABELA EMENTA 23: Filosofia — Ano 1
% ---------------------------------------------------------
\begin{xltabular}{\linewidth}{|X|p{2.5cm}|p{2.5cm}|}
\hline
\multirow{3}{=}{\textcolor{ifscgreen}{\textbf{Unidade Curricular:}}\\[3pt]\textbf{\large Filosofia — Ano 1}} & \multicolumn{2}{p{\dimexpr5cm+2\tabcolsep+\arrayrulewidth\relax}|}{\textcolor{ifscgreen}{\textbf{Semestre:}} \textbf{1º e 2º (Filosofia I)}} \\ \cline{2-3}
 & \textcolor{ifscgreen}{\textbf{CH EaD*:}} & \textcolor{ifscgreen}{\textbf{CH Total*:}} \\ \cline{2-3}
 & \textbf{00 h} & \textbf{80 h} \\ \hline
\endhead
\multicolumn{3}{|p{\dimexpr\linewidth-2\tabcolsep-2\arrayrulewidth\relax}|}{\textcolor{ifscgreen}{\textbf{Objetivos:}}} \\ \hline
\multicolumn{3}{|p{\dimexpr\linewidth-2\tabcolsep-2\arrayrulewidth\relax}|}{
\begin{itemize}[leftmargin=*,noitemsep,topsep=2pt]
 \item Conhecer a origem da Filosofia, os campos de investigação da Filosofia e os períodos da história da Filosofia.
 \item Estimular os primeiros contatos com a tradição filosófica e com a abstração conceitual típica da filosofia.
 \item Adquirir conhecimentos sobre a estrutura do pensamento metafísico através da história da Filosofia.
 \item Conhecer aspectos gerais da filosofia dos pré-socráticos, Sócrates e Platão.
 \item Conhecer aspectos gerais da filosofia de Aristóteles.
 \item Conhecer as escolas filosóficas inerentes ao período helenístico.
 \item Relacionar os conteúdos analisados em sala de aula aos mais distintos aspectos da vida contemporânea.
 \item Promover a troca de ideias e, por conseguinte, o respeito às distintas visões de mundo presentes na sala de aula.
\end{itemize}
} \\ \hline
\multicolumn{3}{|p{\dimexpr\linewidth-2\tabcolsep-2\arrayrulewidth\relax}|}{\textcolor{ifscgreen}{\textbf{Conteúdos:}}} \\ \hline
\multicolumn{3}{|p{\dimexpr\linewidth-2\tabcolsep-2\arrayrulewidth\relax}|}{
\begin{itemize}[leftmargin=*,noitemsep,topsep=2pt]
 \item O que é o Mito? Passagem do Mito à Filosofia. Conceitos de Filosofia: o que é Filosofia e para que serve. Filosofia e as outras formas de conhecimento. A Razão: regras e princípios. As concepções de verdade — Dogmatismo e busca da verdade. Períodos da História da Filosofia. Filosofia Antiga. A Filosofia grega e os pré-socráticos. Os sofistas. Sócrates e a busca pela verdade. Platão e o mito da caverna. O amor a partir de Platão. As cavernas contemporâneas. Aristóteles, ética e política. O homem como animal social. As virtudes e a moderação. A prudência. A coragem. As escolas helenísticas. O estoicismo e a arte de bem viver. O epicurismo e o prazer. O ceticismo e o conhecimento. O cinismo e o cosmopolitismo.
\end{itemize}
} \\ \hline
\multicolumn{3}{|p{\dimexpr\linewidth-2\tabcolsep-2\arrayrulewidth\relax}|}{\textcolor{ifscgreen}{\textbf{Estratégias de Ensino e Aprendizagem:}}} \\ \hline
\multicolumn{3}{|p{\dimexpr\linewidth-2\tabcolsep-2\arrayrulewidth\relax}|}{- **Metodologia:** Aulas expositivas, leituras interpretativas e críticas, seminários e apresentações individuais e/ou em grupo, pesquisa bibliográfica e produção de textos.
- **Avaliação:** Processual, com notas de 0 a 10, considerando participação dos alunos nas atividades propostas, trabalho de pesquisa individual e coletiva e prova escrita, com dois momentos avaliativos formais durante o semestre.} \\ \hline
\multicolumn{3}{|p{\dimexpr\linewidth-2\tabcolsep-2\arrayrulewidth\relax}|}{\textcolor{ifscgreen}{\textbf{Bibliografia Básica:}}} \\ \hline
\multicolumn{3}{|p{\dimexpr\linewidth-2\tabcolsep-2\arrayrulewidth\relax}|}{1. ARANHA, M. L. A.; MARTINS, M. H. P. \textbf{Filosofando: introdução à filosofia}. São Paulo: Moderna, 1993.
2. CHAUÍ, Marilena. \textbf{Convite à filosofia}. 13. ed. São Paulo: Ática, 2009.} \\ \hline
\multicolumn{3}{|p{\dimexpr\linewidth-2\tabcolsep-2\arrayrulewidth\relax}|}{\textcolor{ifscgreen}{\textbf{Bibliografia Complementar:}}} \\ \hline
\multicolumn{3}{|p{\dimexpr\linewidth-2\tabcolsep-2\arrayrulewidth\relax}|}{1. GAARDER, Jostein. \textbf{O mundo de Sofia: romance da história da filosofia}. São Paulo: Seguinte, 2012.
2. MARCONDES, Danilo. \textbf{Textos básicos de filosofia: dos pré-socráticos a Wittgenstein}. Rio de Janeiro: Jorge Zahar, 2000.} \\ \hline
\end{xltabular}

\vspace{0.4cm}


% ---------------------------------------------------------
% TABELA EMENTA 24: Filosofia — Ano 3
% ---------------------------------------------------------
\begin{xltabular}{\linewidth}{|X|p{2.5cm}|p{2.5cm}|}
\hline
\multirow{3}{=}{\textcolor{ifscgreen}{\textbf{Unidade Curricular:}}\\[3pt]\textbf{\large Filosofia — Ano 3}} & \multicolumn{2}{p{\dimexpr5cm+2\tabcolsep+\arrayrulewidth\relax}|}{\textcolor{ifscgreen}{\textbf{Semestre:}} \textbf{5º e 6º (Filosofia 2)}} \\ \cline{2-3}
 & \textcolor{ifscgreen}{\textbf{CH EaD*:}} & \textcolor{ifscgreen}{\textbf{CH Total*:}} \\ \cline{2-3}
 & \textbf{00 h} & \textbf{80 h} \\ \hline
\endhead
\multicolumn{3}{|p{\dimexpr\linewidth-2\tabcolsep-2\arrayrulewidth\relax}|}{\textcolor{ifscgreen}{\textbf{Objetivos:}}} \\ \hline
\multicolumn{3}{|p{\dimexpr\linewidth-2\tabcolsep-2\arrayrulewidth\relax}|}{
\begin{itemize}[leftmargin=*,noitemsep,topsep=2pt]
 \item Conhecer os aspectos gerais da filosofia medieval cristã.
 \item Refletir acerca de autores e questões atinentes à filosofia moderna.
 \item Conhecer o conceito de subjetividade a partir de Montaigne.
 \item Compreender a centralidade da atividade política a partir de Maquiavel, Hobbes, Rousseau e Locke.
 \item Estudar a ética do dever de Kant, bem como a centralidade do conceito de esclarecimento no período iluminista.
 \item Conhecer aspectos gerais do irracionalismo e a crise da razão.
 \item Refletir acerca da crise de valores morais e a problemática do niilismo.
 \item Compreender o existencialismo e os conceitos de projeto e liberdade, investigando a relação entre existencialismo e feminismo.
 \item Visualizar o trabalho como elemento transformador na vida do homem.
 \item Refletir acerca de questões éticas relacionadas à relação homem e meio ambiente.
 \item Promover a troca de ideias e o respeito às distintas visões de mundo presentes na sala de aula.
\end{itemize}
} \\ \hline
\multicolumn{3}{|p{\dimexpr\linewidth-2\tabcolsep-2\arrayrulewidth\relax}|}{\textcolor{ifscgreen}{\textbf{Conteúdos:}}} \\ \hline
\multicolumn{3}{|p{\dimexpr\linewidth-2\tabcolsep-2\arrayrulewidth\relax}|}{
\begin{itemize}[leftmargin=*,noitemsep,topsep=2pt]
 \item Filosofia Medieval Cristã. A relação entre fé e razão. A Patrística de Santo Agostinho. A Escolástica de São Tomás de Aquino. A Filosofia Moderna. A filosofia política em Maquiavel. O pensamento ensaístico de Montaigne. Locke, Rousseau e Hobbes: a filosofia política. Filosofia moderna e contemporânea. Kant, a ética do dever e o esclarecimento. Schopenhauer, pessimismo, vontade e irracionalismo. Nietzsche e a crise dos valores. Niilismo e a crise da razão. Existencialismo e liberdade. A condição da mulher nos séculos XX e XXI. Existencialismo, feminismo e liberdade. Política, Estado e Poder na contemporaneidade. Ética e meio ambiente. O mundo do trabalho e a alienação. Indústria Cultural e consumo.
\end{itemize}
} \\ \hline
\multicolumn{3}{|p{\dimexpr\linewidth-2\tabcolsep-2\arrayrulewidth\relax}|}{\textcolor{ifscgreen}{\textbf{Estratégias de Ensino e Aprendizagem:}}} \\ \hline
\multicolumn{3}{|p{\dimexpr\linewidth-2\tabcolsep-2\arrayrulewidth\relax}|}{- **Metodologia:** Aulas expositivas, leituras interpretativas e críticas, seminários e apresentações individuais e/ou em grupo, pesquisa bibliográfica e produção de textos.
- **Avaliação:** Processual, com notas de 0 a 10, considerando participação dos alunos nas atividades propostas, trabalho de pesquisa individual e coletiva e prova escrita, com dois momentos avaliativos formais durante o semestre.} \\ \hline
\multicolumn{3}{|p{\dimexpr\linewidth-2\tabcolsep-2\arrayrulewidth\relax}|}{\textcolor{ifscgreen}{\textbf{Bibliografia Básica:}}} \\ \hline
\multicolumn{3}{|p{\dimexpr\linewidth-2\tabcolsep-2\arrayrulewidth\relax}|}{1. ARANHA, M. L. A.; MARTINS, M. H. P. \textbf{Filosofando: introdução à filosofia}. São Paulo: Moderna, 1993.
2. CHAUÍ, Marilena. \textbf{Convite à filosofia}. 13. ed. São Paulo: Ática, 2009.} \\ \hline
\multicolumn{3}{|p{\dimexpr\linewidth-2\tabcolsep-2\arrayrulewidth\relax}|}{\textcolor{ifscgreen}{\textbf{Bibliografia Complementar:}}} \\ \hline
\multicolumn{3}{|p{\dimexpr\linewidth-2\tabcolsep-2\arrayrulewidth\relax}|}{1. GAARDER, Jostein. \textbf{O mundo de Sofia: romance da história da filosofia}. São Paulo: Seguinte, 2012.
2. MARCONDES, Danilo. \textbf{Textos básicos de filosofia: dos pré-socráticos a Wittgenstein}. Rio de Janeiro: Jorge Zahar, 2000.} \\ \hline
\end{xltabular}

\vspace{0.4cm}

